\documentclass[11pt]{article}
\usepackage[a4paper, top=18mm, bottom=18mm, left=20mm, right=20mm]{geometry}
\usepackage[T2A]{fontenc}
\usepackage[utf8]{inputenc}
\usepackage[ukrainian]{babel}
\usepackage{graphicx}
\usepackage[export]{adjustbox}
\usepackage{amsmath}
\usepackage{amssymb}
\setlength{\parindent}{0pt}
\setlength{\parskip}{4pt}
\pagestyle{empty}
\begin{document}
%begin
{\large\textbf{Контрольна робота}}

\textbf{Варіант \No\,1}\hfill \textit{ПІБ:}\,\underline{\hspace{60mm}}
\vspace{4mm}

%beginitem
1. %goodpath:Scripts/2018/Mod2018_1/03-good.txt
Відмітити все, що є коректним оператором С++:
( 1 ) double n.m;

( 2 ) if a<b a=0;

( 3 ) {int x,y;}

( 4 ) double(2)/87

%enditem

%beginitem
2. 
try:
    print(5, end=';')
    print(2j!=5j, end=';')
    print(0, end=';')
except TypeError:
    print(9, end=';')
except ZeroDivisionError:
    print(6, end=';')
except:
    print('ERR', end=';')
else:
    print(8, end=';')
finally:
    print(1)

%enditem

%beginitem
3. 
code
d = 0
k = 1

class D:
    k = 2
    
    def __init__(self):
global k        self.k = 3
        k = 4
        D.k = 5

obj = D()
print(d, k, D.k, obj.k, sep=' ')
code

%enditem

%beginitem
4. Що буде виведено за виконання фрагменту коду?
%code
if (5 != 5)
    cout << "f";
else
    cout << "s";
%code


%enditem

%beginitem
5. 
Записати ВИРАЗ, значенням якого є словник, ключами якого є цілі числа від 12 до 2012 (включно), що не діляться на жодне з чисел 2, 3, а ключам відповідають їх квадратні корені 
%enditem

%beginitem
6. %goodpath:Scripts/2019/Mod2019_1/Lex/01-good.txt
Відмітити все, що є однією правильною лексемою:
( 1 ) 5_w

( 2 ) 0x17

( 3 ) 'c'

( 4 ) "1,3,2"

%enditem

%beginitem
7. Що буде виведено за виконання фрагменту коду?

%code
a,b,c=5,2,0
def h(b):
global c    a*=3
    b=4
    c=2
    return a+b+c

a,b,c=9,6,8
print(h(b),a,b,c)
%code

%enditem

%beginitem
8. 
try:
    res = 9**-1.0*False
    if isinstance(res, bool):
        print(f'bool;{res}')
    elif isinstance(res, int):
        print(f'int;{res}')
    elif isinstance(res, float):
        print(f'float;{res:.2f}')
    else:
        print('???')
except:
    print('Z')

%enditem

%beginitem
9. Що буде виведено за виконання фрагменту коду?

%code
def h(a,b):
    c=84
    if b!=-3:
        c=5
    elif a>-2:
         return 2
    else: 
        return 0
    return c

print(h(3,3))
%code

%enditem

%beginitem
10. %goodpath:Scripts/2019/Mod2019_1/Lex/03-good.txt
Відмітити все, що є коректним оператором С++:
( 1 ) x not is y

( 2 ) print(end="1","qwe")

( 3 ) 3**-1

( 4 ) x-=2

%enditem

%end
\newpage
\end{document}

